\section{Support acquisition at the fifth index}\label{ext:n5-support}
The finite Fourier and count estimates require an explicit support radius.
Compactness alone is insufficient: the radius enters their error bounds.
Here we obtain the radii in the following table directly from contact at
$n=5$, uniformly in the raw threshold. Corollary~\ref{ext:n5-refined-support} gives the tighter radii used by the
signed-moment sectors of the initial acquisitions.

\begin{proposition}[The fifth-index support table]\label{ext:n5-support-table}
Assume $C_4\le\CE$ and $C_5>\CE$. Let $(P,H)$ be a positive
closed-event ratio maximizer supplied by Corollary~\ref{ext:ratio-attainment}.
If $\rho(P)$ belongs to a row in the following table, then
$\supp P\subset(-R,R)$ for that row's radius. The values of $b$
are upper bounds for $|\E[X|X|]|$. The last column is a rational
lower bound on the strict endpoint margin in the proof.
\begin{center}\small
\begin{tabular}{@{}cccc@{}}
\toprule
$\rho$ interval & $b$ & $R$ & endpoint margin $>$\\
\midrule
$[1,11/10]$ & $0.253257$ & $23/10$ & $0.251666$\\
$[11/10,23/20]$ & $0.307221$ & $12/5$ & $0.153351$\\
$[23/20,237/200]$ & $0.338920$ & $5/2$ & $0.174691$\\
$[237/200,6/5]$ & $0.351391$ & $13/5$ & $0.350224$\\
$[6/5,5/4]$ & $0.389174$ & $13/5$ & $0.020061$\\
$[5/4,7/5]$ & $0.478710$ & $117/40$ & $0.005870$\\
$[7/5,3/2]$ & $0.525521$ & $25/8$ & $0.008946$\\
$[3/2,8/5]$ & $0.565414$ & $133/40$ & $0.059409$\\
$[8/5,17/10]$ & $0.600001$ & $7/2$ & $0.033467$\\
$[17/10,7361/4000]$ & $0.641563$ & $15/4$ & $0.061370$\\
\bottomrule
\end{tabular}
\end{center}
\end{proposition}

\begin{proof}
Write $C=C_5$, $\beta=\E[X|X|]$, $s=\sqrt{5/4}$ and
$e=5(s-1)-1/2$. The score of the same marginal at index four is
$A=2\Delta_4(P)/\rho\le C_4<C$. Hence the signed predecessor term
$d_P=5(A-C)$ is nonpositive.
Define
\[
 \mathcal H(H,t)=5\sqrt5\{\Phi((H-t)/2)-\Phi(H/\sqrt5)\}
 +5t\phi(H/\sqrt5)+\tfrac H2(t^2-1)\phi(H/\sqrt5).
\]
Contact and the predecessor bound imply, at every support point $t$,
\begin{equation}\label{ext:n5-contact-envelope}
 0\le\mathcal H(H,t)
 +C\{-|t|^3+\tfrac32\rho t^2+3\beta t+\rho e\}
 +\rho s d_P.
\end{equation}
To verify the scale term, the probability part contributes
$5\rho(sA-C)$. The identity
$5(sA-C)=s d_P+5(s-1)C$ combines with the constant
$-C\rho/2$ in the contact polynomial to give exactly $C\rho e$.
Thus there is no approximation in \eqref{ext:n5-contact-envelope}.

For a row with upper moment endpoint $r$, put
\[
 c=819/2000,\qquad E=91/1000,\qquad
 Q(v)=-v^3+\tfrac32rv^2+3bv+rE,
\]
\[
 L(v)=\tfrac25(v^2/2+1/2+E)+cQ'(v).
\]
The elementary constants satisfy $c<\CE<C$, $e<E$ and
$s<1119/1000$. The signed-moment inequality of Lemma~\ref{ms:sign-moment},
$(1+\beta^2)^2\le\rho^2(1-\beta^2)$ implies $|\beta|\le b$,
because direct rational substitution gives
$(1+b^2)^2>r^2(1-b^2)$ for each row.
The exact arithmetic check gives
\begin{equation}\label{ext:n5-ray-gates}
 Q(R),Q'(R),Q''(R),L(R),L'(R)<0,
 \qquad 1/5-3c<0.
\end{equation}
Since $Q'''=-6$, these gates imply $Q(v)<0$ for all $v\ge R$.
The bracket multiplying $C$ in \eqref{ext:n5-contact-envelope} is
at most $Q(|t|)$, so on those rays its product with $C$ is at most
$cQ(|t|)$. The predecessor term is nonpositive.

The Gaussian envelope has the uniform derivative bound
\begin{equation}\label{ext:n5-ray-derivative}
 |\partial_t\mathcal H(H,t)|
 \le\phi(0)(t^2/2+1/2+e)
 <\tfrac25(t^2/2+1/2+E).
\end{equation}
Indeed, with $z=H/\sqrt5$ and $y=(H-t)/\sqrt5$,
$\partial_t\mathcal H=5\phi(z)+Ht\phi(z)-5s\phi(sy)$.
Taylor's theorem and $\|\phi''\|_\infty=\phi(0)$ bound
$|5\phi(y)-5\phi(z)-Ht\phi(z)|$ by $\phi(0)t^2/2$.
Also $\partial_s[s\phi(sy)]=-\phi''(sy)$, so the remaining
error is at most $5\phi(0)(s-1)$. This proves
\eqref{ext:n5-ray-derivative} for every real $H,t$.
For either sign $\eta$, the derivative in $v\ge R$ of
$\mathcal H(H,\eta v)+cQ(v)$ is at most $L(v)$.
Now $L''=2/5-6c<0$; the last two gates in
\eqref{ext:n5-ray-gates} make $L$ strictly negative throughout the ray.
It remains to prove a strict upper bound at $v=R$, on the whole
threshold line.

For that purpose, Gaussian differentiation gives
\[
 |\partial_H^2\mathcal H(H,R)|
 \le\frac{\sqrt5}{4}\left(\frac54+1\right)
       +\frac25R+\frac{2|R^2-1|}{5\sqrt5}
 \le M_R:=\frac{81}{64}+\frac25R+\frac{|R^2-1|}{5}.
\]
The first bound uses $\|\phi'\|_\infty=1/\sqrt{2\pi e}<1/4$
(the maximum of $|x|e^{-x^2/2}$ occurs at $|x|=1$),
$\|\phi''\|_\infty<2/5$ and
$\|\phi'''\|_\infty\le2/\pi<4/5$; the last follows either
by differentiation or by the Gaussian Fourier integral. The second
bound uses $2<\sqrt5<9/4$.
For a twice differentiable function with $|f''|\le M$ on $[a,b]$,
its deviation from the affine endpoint interpolant is at most
$M(b-a)^2/8$. This follows by comparing the interpolation error
with $\pm M(x-a)(b-x)/2$. Consequently, the values at
\[
 H=-16,-16+1/10,\ldots,16
\]
plus $M_R/800$ bound $|\mathcal H(H,R)|$ on every one of the
320 closed intervening cells.

The two exterior rays are covered separately. Since $16\ge5\sqrt5$
and $16-R\ge10$, Mills' bound gives, with $\tau=10^{-6}>\Phi(-5)$,
\[
 |\mathcal H(H,R)|\le5\sqrt5\,\tau+
 \left(5R+8|R^2-1|\right)\phi(16/\sqrt5),\qquad |H|\ge16.
\]
Here $|H|\phi(H/\sqrt5)$ decreases for $|H|\ge16$.
Finally $\mathcal H(H,-R)=-\mathcal H(-H,R)$, so the same
absolute enclosure handles both support rays.
The resulting whole-line bound $B_R$ satisfies
$B_R+cQ(R)<0$; the table records lower bounds for $-cQ(R)-B_R$.
Together with the ray monotonicity, this contradicts
\eqref{ext:n5-contact-envelope} whenever $|t|\ge R$.
\end{proof}

\paragraph{Exact evaluation of the finite table.}
The check is a short independent rational program, rather than a sampled
numerical maximum. It brackets $\pi$ using the alternating arctangent
series and $\pi=16\arctan(1/5)-4\arctan(1/239)$, and brackets positive
square roots by integer squaring on the grid $10^{-25}$.
For $0\le x\le26$, truncating $e^{-x}$ after the even term 120 and
adding the negative next term gives an enclosing interval. For
$H^2<25v$, the formula
\[
 \Phi(H/\sqrt v)=\frac12+
 \frac{\phi(0)}{\sqrt v}
 \sum_{k=0}^{\infty}
 \frac{(-1)^kH^{2k+1}}{(2v)^k k!(2k+1)}
\]
is truncated after the even term 80, with the next term as remainder
bound. The relevant tails of both alternating series decrease in absolute
value. For $H^2\ge25v$, use $\Phi(-5)<10^{-6}$ and reflection.
All interval operations retain their signs. The check evaluates all 3,200
closed threshold cells in the ten rows, both exterior rays, every
polynomial gate, the signed-moment enclosures and the constants.
Insufficient support, an invalid moment bound and a threshold mesh too
coarse to prove the assertion are rejected in separate controls.
The full fractions are saved with the execution record; displayed decimals
are exact rationals rounded down from the proved margins.

\begin{corollary}[Radii for the fifth-index Fourier and resource bounds]
\label{ext:n5-refined-support}
Under the assumptions of Proposition~\ref{ext:n5-support-table}, one has
\[
 \begin{array}{c|c|c|c}
 \rho\text{ interval}&|\beta|\text{ upper bound}&R&
       \text{endpoint margin}>\\ \hline
 {[1,11/10]}&0.253257&11/5&0.045938\\
 {[11/10,23/20]}&0.307221&47/20&0.039564\\
 {[23/20,237/200]}&0.338920&49/20&0.050250\\
 {[237/200,6/5]}&0.351391&99/40&0.016341
 \end{array}
\]
In each row, $\supp P\subset(-R,R)$. These radii hold on the
whole signed-moment range and therefore on every narrower signed sector
used by the resource calculation.
\end{corollary}
\begin{proof}
Use exactly the polynomial gates, interpolation bound, Gaussian enclosures,
and two infinite tails from Proposition~\ref{ext:n5-support-table}.
The additional exact checks use the four displayed radii,
on the same 320 closed threshold cells per row. All five ray gates and
the strict endpoint inequalities pass. The lower margins shown in the
table are rounded downward from the exact rational margins. The trial radii $21/10$, $9/4$, $12/5$ and $12/5$, respectively,
fail the acceptance criterion. These failures are rejection controls,
not counterexamples to a support theorem.
\end{proof}
